












\DeclareFontFamily{U}{mathx}{\hyphenchar\font45}
\DeclareFontShape{U}{mathx}{m}{n}{
  <5> <6> <7> <8> <9> <10>
  <10.95> <12> <14.4> <17.28> <20.744> <24.888>
  mathx10
}{}
\DeclareSymbolFont{mathx}{U}{mathx}{m}{n}
\DeclareMathAccent{\widecheck}{0}{mathx}{"71}

\newcommand{\somme}{\sum} %Displaystyle sum
\newcommand{\produ}{\prod} %Displaystyle product
\newcommand{\tsp}[1]{{}^t \! #1} %French transpose
\newcommand{\ii}{\mathrm{i}} %Imaginary Unit
\newcommand{\tr}{\operatorname{Tr}} %Trace
\newcommand{\Span}{\operatorname{Span}} %Linear span
\newcommand{\supp}{\operatorname{supp}} %Support
\newcommand{\im}{\operatorname{im}} %Image
\newcommand*\diff{\mathop{}\!\mathrm{d}} %Differential
\newcommand{\Diff}{\operatorname{Diff}} %Diffeomorphisms or differential operators
\newcommand{\ssing}{\operatorname{ssing}} %Singular support
\newcommand{\Id}{\operatorname{Id}} %Identity
\newcommand{\R}{\mathbb R} %Real numbers
\newcommand{\N}{\mathbb N} %Natural numbers
\newcommand{\Q}{\mathbb Q} %Rational numbers
\newcommand{\Z}{\mathbb Z} %Integers
\newcommand{\CC}{\mathbb C} %Complex numbers
\newcommand{\A}{\mathcal{A}} %Lie algebroid
\newcommand{\G}{\mathcal{G}} %Fancy groupoid
\newcommand{\hol}{\operatorname{Hol}} %Holonomy
\newcommand{\CCC}{\mathscr{C}} %Continuous functions
\newcommand{\F}{\mathcal{F}} %Foliation
\newcommand{\PP}{\mathbb{P}} %Big bold P
\newcommand{\card}{\operatorname{Card}} %Cardinal
\newcommand{\GL}{\operatorname{GL}} %General linear group
\newcommand{\Hom}{\operatorname{Hom}} %Homomorphisms
\newcommand{\End}{\operatorname{End}} %Endomorphism algebra
\newcommand{\Aut}{\operatorname{Aut}} %Automorphism group
\newcommand{\Ad}{\operatorname{Ad}} %Adjoint group action
\newcommand{\ad}{\operatorname{ad}} %Adjoint Lie algebra action
\newcommand{\codim}{\operatorname{codim}} %Codimension
\newcommand{\zz}{\mathbf{z}} %Bold z
\newcommand{\g}{\mathfrak{g}} %Lie algebra
\newcommand{\p}{\mathfrak{p}} %Parabolic Lie subalgebra
\newcommand{\n}{\mathfrak{n}} %Nilpotent Lie algebra
\newcommand{\adiab}{\mathfrak{a}} %Adiabatic
\newcommand{\Heis}{\operatorname{Heis}} %Heisenberg group
\newcommand{\heis}{\mathfrak{heis}} %Heisenberg Lie algebra
\newcommand{\delbar}{\bar{\partial}}%Delbar operator
\newcommand{\vphi}{\varphi}
\newcommand{\delbarE}{\bar{\partial}_E} %Delbar operator for bundle
\newcommand{\dn}{\diff^{\nabla}} %Connection differential
\newcommand{\ind}{\operatorname{Ind}} %Index
\newcommand{\partialt}{\frac{\partial}{\partial t}} %Partial time differential
\newcommand{\rr}{\rightrightarrows} %Groupoid
\newcommand{\surj}{\twoheadrightarrow} %Map is a surjection
\newcommand{\inj}{\hookrightarrow} %Map is an injection
\newcommand{\act}{\curvearrowright} %Action
\newcommand{\isom}{\xrightarrow{\sim}} %Map is an isomorphism
\newcommand{\lag}{\langle}
\newcommand{\rag}{\rangle}
\newcommand{\DNC}{\operatorname{DNC}} %Deformation to the normal cone
\newcommand{\op}{\operatorname{Op}} %Associated operator
\newcommand{\ttt}{\mathfrak{t}} %t lie algebra
\newcommand{\T}{\mathcal{T}} %Toeplitz algebra, fancy curly T
\newcommand{\Lie}{\operatorname{Lie}} %Lie functor
\newcommand{\hd}{\Omega^{1/2}} %Half-densities
\newcommand{\rk}{\operatorname{rk}} %Rank
\newcommand{\m}{\operatorname{\mathcal{M}}} %Multiplier
\newcommand{\pr}{\operatorname{pr}} %Projection
\newcommand{\ev}{\operatorname{ev}} %Evaluation map
\newcommand{\ord}{\operatorname{ord}} %Order
\newcommand{\Sym}{\operatorname{Sym}} %Symmetric tensor
\newcommand{\HN}{\operatorname{HN}} %Helffer--Nourrigat
\newcommand{\Fb}{\mathcal{F}} %Singular filtration
\newcommand{\gr}{\mathfrak{gr}} %For associated grading
\newcommand{\U}{\mathbb U} %For the graded basis

\newcommand\isomto{\stackrel{\sim}{\smash{\longrightarrow}\rule{0pt}{0.4ex}}}

\newcommand*{\corner}{\mbox{\LARGE{$\mathrlap{\cdot}\lrcorner$}}}
\def\dar[#1]{\ar@<2pt>[#1]\ar@<-2pt>[#1]}
\entrymodifiers={!!<0pt,0.7ex>+}

\newcommand{\eq}[1][r]
 {\ar@<-3pt>@{-}[#1]
 \ar@<-1pt>@{}[#1]|<{}="gauche"
 \ar@<+0pt>@{}[#1]|-{}="milieu"
 \ar@<+1pt>@{}[#1]|>{}="droite"
 \ar@/^2pt/@{-}"gauche";"milieu"
 \ar@/_2pt/@{-}"milieu";"droite"}
\everymath{%\displaystyle
\everymath{}}

\numberwithin{equation}{section}

\newtheorem{thm}{Theorem}[section]
\newtheorem{propr}[thm]{Property}
\newtheorem{lem}[thm]{Lemma}
\newtheorem{prop}[thm]{Proposition}
\newtheorem{cor}[thm]{Corollary}

\theoremstyle{definition}
\newtheorem{mydef}[thm]{Definition}
\newtheorem{ex}[thm]{Example}
\newtheorem{rem}[thm]{Remark}
\newtheorem*{nb}{NB}
































\section{Introduction}

Let $X_1, \dots, X_m$ be vector fields on a~manifold $M$ that are Lie bracket generating. This means that for any $x \in M$, the linear span of vectors of the form $X_I(x), I\subset \{1;\dots;m\}^{\N}$ is equal to~$T_xM$, where
\[X_{i_1,\dots,i_N} := [X_{i_1},[\dots,[X_{i_{n-1}},X_{i_N}]\dots].
\]
This condition, known as H\"ormander condition, appeared in~\cite{Hormander}, where it was proved to be a~sufficient condition for the hypoellipticity of the sublaplacian $\sum_{i=1}^m X_i^2$. A~differential operator~$P$ is hypoelliptic if for any distribution $u \in \CCC^{-\infty}(M)$ such that $Pu$ is smooth then $u$ is also smooth.
Given such vector fields, a~natural question is how to determine the hypoellipticity of a~polynomial in the vector fields $X_1, \dots, X_m$, beyond the sublaplacian (i.e., a~sum of squares). A~necessary and sufficient condition for maximal hypoellipticity (a stronger version of hypoellipticity mimicking the usual elliptic estimates) was conjectured in~\cite{HelfferNourrigatPaper} and solved completely in the recent work~\cite{AMYMaximalHypoellipticity}, see~\cite{DebordBourbaki,HelfferNourrigatBook} for a~more complete overview and history of the problem.

To solve this conjecture, the authors of~\cite{AMYMaximalHypoellipticity} develop a~pseudodifferential calculus, and show that the invertibility criterion of the principal symbol (i.e., the ellipticity criterion in this calculus) is exactly what was conjectured by Helffer and Nourrigat. The idea of this new calculus is to consider the vector fields $X_{i_1,\dots,i_N}$ as having order $N$. More precisely, a~differential operator that can be written as a~$\CCC^{\infty}(M)$-linear combination of the vector fields $X_{I}$, $|I|\leq N$ has order at most $N$, its order being the minimal such natural number $N$. Because of the H\"ormander condition, every vector field can be written in terms of the vector fields $X_I$, so every differential operator has a~finite, well-defined order. In order to define a~principal symbol at a~point $x\in M$, the idea is to freeze the coefficients at $x$ and look at a~model operator. In the usual calculus, this operator is a~differential operator with constant coefficients on the tangent space $T_xM$. Here, because of the inhomogeneity in the directions, the model operator becomes a~differential operator on a~graded nilpotent group arising naturally from the anisotropy: the osculating groups~\cite{AMYMaximalHypoellipticity,RothschildStein}. In such a~general setting, these nilpotent groups can have a~dimension bigger than the dimension of $M$, they will typically be universal nilpotent groups following the methods of~\cite{RothschildStein}. To understand the model operator, one uses its image in the unitary representations of the group. Not all these representations make sense for our initial operator on $M$, giving rise to the Helffer--Nourrigat cone. This cone therefore plays the role of a~phase space for the calculus at hand as it replaces the cotangent bundle. Indeed, in this sub-elliptic setting, certain phenomena such as observability and controllability (for associated wave or Schr\"odinger operators) are not captured by the classical phase space, but by this singular ``non-commutative'' phase space, see, e.g., \cite{FermanianLetrouit,Letrouit}

The unitary representations of a~nilpotent group form a~topological space with the Fell topology. It was proved by Kirillov~\cite{Kirillov} that this space is homeomorphic to the space of coadjoint orbits, i.e., the dual of the Lie algebra, quotiented by the coadjoint action and endowed with the quotient topology. As such, this space can be quite pathological, usually being not Hausdorff. We can however stratify this space into some nice parts following Puk\'anszky and Pedersen. On these parts the representations can be continuously parameterized, in particular the corresponding subquotient of the group $C^*$-algebra is an algebra of continuous functions (with values in compact operators). In the case of graded groups, this stratification can be also applied to the spectrum of the algebra of principal symbols on the group. It was shown in~\cite{CrenFilteredCrossedProduct} that one can do this stratification over the whole manifold $M$ when the filtration is equiregular, see Section~\ref{Section:PedersenStratification} below. The strata appearing in the Pedersen decomposition of the space of representations of a~nilpotent group can be identified to semi-algebraic sets in the vector space $\g^*$. Therefore, one can hope to get a~similar stratification for the Helffer--Nourrigat cone and use it to adapt geometric constructions used in the classical setting. For instance, can the notion of Hamiltonian dynamics on the cotangent bundle $T^*M$ make sense on the Helffer--Nourrigat cone? If so, one can hope to understand questions such as propagation of singularities, of semiclassical measures, define a~geometric control condition, study quantum ergodicity (see, e.g., \cite{benedetto}). Such stratification can also be used in index theory as in~\cite{GoffengKuzmin}.

We wish here to extend this stratification in a~general setting. More precisely, given a~family of vector fields on $M$ satisfying H\"ormander's condition, we denote by $\gr(\F)$ the family of osculating groups over $M$ ($\F$ denotes the singular Lie filtration of the tangent bundle induced by the vector fields, see below). We denote by \smash{$\HN(\F) \subset \widehat{\gr(\F)}$} the Helffer--Nourrigat cone, it is a~closed subset of the irreducible unitary representations of all the osculating groups. For $x \in M$, the fiber of the cone at $x$ contains the trivial representation of the osculating group $\gr(\F)_x$. We consider $M$ as a~subset of $\HN(\F)$ through the trivial representations of the fibers and denote by $\HN_0(\F)$ the complement of this inclusion. The spaces $\HN( \F)$ and $\HN_0(\F)$ are endowed with an action of~$\R^*_+$ obtained from the inhomogeneous dilations of the underlying osculating groups. This action is free and proper on $\HN_0(\F)$. We prove the following theorem.

\begin{thm}[stratification of the Helffer--Nourrigat cone]\label{Theorem:Stratification}
 Let $\Fb$ be a~finitely generated singular Lie filtration on a~manifold $M$. There exists a~filtration by open sets
 \[\varnothing = U_0 \subset U_1 \subset \cdots \subset U_d = \HN(\Fb),\]
 such that all the spaces $\Omega_j = U_j\setminus U_{j-1}$ are locally compact Hausdorff.

 We also have \smash{$\HN_0(\Fb) = \bigsqcup_{j=1}^{d-1} \Omega_j \sqcup (\Omega_d \setminus M)$}.
 The action of $\R^*_+$ on $\HN(\Fb)$ induced by the inhomogeneous dilation on $\gr(\Fb)$ preserves these strata. The spaces \smash{${\Omega_j}/{\R^*_+}$}, $1\leq j < d$ and~\smash{${(\Omega_d\setminus M)}/{\R^*_+}$} are locally compact Hausdorff as well.

 Finally, for $x\in M$, denote by $\Omega_{j,x} = \Omega_j\cap \HN(\Fb)_{x}$. Then, if $x \in M_{\rm reg}$, the non-empty sets in the collection $(\Omega_{j,x})_{1\leq j \leq d}$ form a~Pedersen stratification of the group $\gr(\Fb)_x$. If $x \in M \setminus M_{\rm reg}$, the sets $\Omega_{j,x}$ are closed subsets of the Pedersen strata of some Pedersen stratification of $\gr(\Fb)_x$.
\end{thm}






















\subsection*{Conventions and notations} All the Lie algebras involved in this paper are nilpotent. Because of this, we will use the Baker--Campbell--Hausdorff formula to identify the Lie algebra with the connected, simply connected group integrating it. We will write both of them using lowercase gothic letters (e.g.,~$\g$). The only exception to that rule are some examples of nilpotent groups in Section~\ref{Section: ExamplesCones}, where the notations try to match existing ones in the literature. The notation $\widehat{\g}$ will refer to the space of irreducible unitary representations of the group $\g$ with its Fell topology. We will denote by $0 \in \widehat{\g}$ the trivial representation.

The letter $\CCC$ will be used to denote the space of continuous functions on a~locally compact, Hausdorff space. The subscript $\CCC_c$ will denote the subspace of compactly supported functions and $\CCC_0$ continuous functions vanishing at infinity. The superscript $\CCC^{\infty}$ will denote smooth functions and will be combined with the aforementioned subscripts.

\section{The Helffer--Nourrigat cone} \label{Section: Helffer--Nourrigat cone}

\subsection{Generalities on distributions and singular Lie filtrations}

In this section, we describe the Helffer--Nourrigat cone of representations. We describe it in two different ways. The first is intrinsic and obtained from the so called collection of osculating groups over each point of the manifold following~\cite{AMYMaximalHypoellipticity}. The second follows the method of Rothschild and Stein~\cite{RothschildStein} and describes the cone as a~subset of representations of a~single graded group (parameterized by the manifold~$M$). This graded group depends on choices of generators for each module $\F^j$, $j\geq 1$, of the filtration. Once these choices are made, they also induce a~map from this group to each osculating group. The dual of this map induces a~homeomorphism between the two constructions of the Helffer--Nourrigat cone.

We start this section by generalities on modules of vector fields. Let $\mathcal{D}\subset \mathfrak{X}(M)$ be a~$\CCC^{\infty}(M)$-sub-module, locally finitely generated. For a~given $x \in M$ we denote by \smash{$\mathcal{D}_x := {\mathcal{D}}/{I_x\mathcal{D}}$} the fiber at $x$, where $I_x= \{ f\in \CCC^{\infty}(M),\, f(x) = 0\}$ is the ideal of functions vanishing at $x$. This vector space is finite-dimensional because $\mathcal{D}$ is locally finitely generated. The evaluation at $x$ induces a~linear map $\mathcal{D}_x \to T_xM$. The image of this map is $D_x = \ev_x(\mathcal{D})$ but the map is not always injective.

The subset $M_{\mathcal{D}} \subset M$ of points over which this map is bijective is called the regular part of $\mathcal{D}$. By~\cite[Proposition~1.5]{AndroulidakisSkandalisHolonomy}, the regular part of $\mathcal{D}$ is an open dense subset of $M$, the family of $D_x, x\in M_{\mathcal{D}}$ forms a~subbundle of $ TM_{\mathcal{D}}$, and the restriction of $\mathcal{D}$ to $M_{|\mathcal{D}}$ is its module of sections. In~\cite{AndroulidakisSkandalisHolonomy}, this is proved for a~singular foliation but the proof of this result only uses the locally finitely generated assumption, so it is also valid in our context.

Finally, denote by $\mathcal{D}^* = \bigsqcup_{x\in M}\mathcal{D}_x^*$. We define a~topology on $\mathcal{D}^*$ by requiring that the projection on $M$ is continuous as well as the maps $\lag \cdot, X \rag, X\in \mathcal{D}$. Following~\cite[Proposition~2.10]{AndroulidakisSkandalisPseudoDiff}, it is a~locally compact Hausdorff space.

\begin{mydef}
 Let $M$ be a~smooth manifold. A~singular Lie filtration is a~family \smash{$\Fb = \bigl(\F^j\bigr){}_{j\geq 1}$} of locally finitely generated $\CCC^{\infty}(M)$-sub-modules of the module of vector fields $\mathfrak{X}(M)$. They satisfy the condition on the Lie brackets $\forall i,j\geq 1$, $\big[\F^i,\F^j\big] \subset \F^{i+j}$.
 Such a~filtration satisfies the H\"ormander condition if $\F^k = \mathfrak{X}(M)$ for $k$ big enough.
\end{mydef}

The H\"ormander condition is technically stronger than the original one. This condition requires that for every $x\in M$, $F^k_x = T_xM$ for $k$ big enough (that might depend on $x$). We can always take a~uniform $k$ locally. Indeed, take $X_1,\dots,X_N \in \F^k$ such that their evaluations~${X_1(x),\dots,X_N(x)}$ at $x\in M$ span $T_xM$. Then there exists an open neighborhood~${x\in U\subset M}$ such that $X_1, \dots, X_N$ is a~basis of $\F^k_{|U}$. In particular, if $M$ is compact, then we can take a~uniform $k$ that works for every point. In the case where the manifold is non-compact, we want to avoid situations where we need more and more depth in the filtration when going to infinity. Therefore, in this paper, we always assume such $k$ to be uniform for all the points in~$M$ and the H\"ormander condition will refer to the definition above, which is slightly stronger than the one in the literature.

\begin{ex}
 Let $X_1, \dots, X_N$ be vector fields on $M$. Set $\F^1 = \CCC^{\infty}(M)\lag X_1,\dots,X_N\rag$, and define inductively
 \[\F^{k+1} = \F^k + \big[\F^k,\F^1\big].\] Then $\F = \bigl(\F^k\bigr)_{k\geq 1}$ is a~singular Lie filtration. The H\"ormander condition on the filtration is then equivalent to the one on the vector fields (from the initial work of H\"ormander~\cite{Hormander}) when $M$ is compact (in general, the weaker version given above is equivalent to the H\"ormander condition on the vector fields).
\end{ex}

\begin{ex}\label{Example : Finitely generated}
 More generally, one can take vector fields $X_1, \dots, X_N$ and associate weights $\nu_1,\dots,\nu_N \in \N$ to them. Then declare that a~vector field $X_I$, $I\in \{1,\dots,N\}^k$, $k\in \N$ has weight~\smash{$\sum_{j\in I}\nu_j$}. Consider then $\F^k$ to be the sub-module of vector fields generated by the ones of weight at most $k$. The family $\F = \bigl(\F^k\bigr)_{k\geq 1}$ forms a~singular Lie filtration on $M$ and the (weaker version of the) H\"ormander condition on $\F$ again corresponds to the one on the vector fields $X_1, \dots, X_N$.

 This singular Lie filtration is (globally) finitely generated. Conversely, every finitely generated singular Lie filtration is of this form. This can be seen by taking a~basis of the first stratum and iteratively adding vector fields not obtainable by Lie brackets to get a~generating set for each~$\F^k$,~$k\geq 1$.
\end{ex}


\subsection{The Helffer--Nourrigat cone: Intrinsic definition}

We now construct the osculating groups associated to a~singular Lie filtration $\Fb$. At a~given point $x\in M$, we define
\[\gr_x(\Fb):= \bigoplus_{j\geq 1} \frac{\F^j_x}{\F^{j-1}_x} = \bigoplus_{j\geq 1} \frac{\F^j}{\F^{j-1} + I_x\F^j},\]
with the convention $\F^0 = \{0\}$.
In the quotient \smash{$\frac{\F^j_x}{\F^{j-1}_x}$}, the map $\F^{j-1}_x \to \F^j_x$ comes from the inclusion between the modules, $\F^{j-1}\subset \F^j$, but is not necessarily injective between the fibers over~$x$. Therefore, the dimension of $\gr_x(\Fb)$ can be greater than the dimension of~$M$.
Given~${i,j\geq 1}$, $x\in M$, and sections $X\in \F^i$, $Y\in \F^j$, the value of $[X,Y](x) \mod \F_x^{i+j-1}$ only depends on~${X(x) \mod \F_x^{i-1}}$ and $Y(x) \mod \F_x^{j-1}$. This defines a~Lie bracket on $\gr_x(\Fb)$, giving it the structure of a~graded nilpotent Lie algebra. Using the Baker--Campbell--Hausdorff formula, this defines a~connected, simply connected, graded Lie group. We can therefore look at the unitary representations of this group, that we understand through Kirillov's orbit method~\cite{Kirillov}
\[\widehat{\gr_x(\Fb)} \cong {\gr_x(\Fb)^*}/{\Ad^*(\gr_x(\Fb))}.\]
To define the Helffer--Nourrigat cone, we need the adiabatic foliation from~\cite{AMYMaximalHypoellipticity}
\[\adiab\F := \sum_{j\geq 1} t^j\widetilde{\F}^j \subset \mathfrak{X}(M\times \R_+).\]
Here, $t$ denotes the projection onto the second factor and $\widetilde{\F}^j$ is the module of vector fields on~${M\times \R_+}$ generated by the elements $(X,0), X \in \F^j$. By the properties of a~singular Lie filtration, this defines a~singular foliation on $M \times \R_+$. For $x\in M$, the point $(x,0)$ is stationary for $\adiab\F$, therefore the fiber $\adiab\F_{(x,0)}$ is a~nilpotent Lie algebra. It is shown in~\cite[Proposition 1.6]{AMYMaximalHypoellipticity} that this Lie algebra is naturally isomorphic to $\gr_x(\Fb)$. For $t \neq 0$, we have $\adiab\F_{(x,t)} = T_xM$ because of the H\"ormander condition.

On the graded Lie algebra $\gr(\F)_x$, there is a~natural $\R^*_+$-action $\delta_{\lambda}, \lambda>0$ given by multiplication by $\lambda^j$ on $\gr(\Fb)_{j,x}$. This is compatible with the Lie bracket and is thus a~Lie algebra automorphism. We denote the Lie group automorphism the same way. One can check that this action is induced from the zoom action defined on $M \times \R_+$ by $\alpha_{\lambda}(x,t) = \bigl(x,\lambda^{-1}t\bigr)$, which preserves the adiabatic foliation.

\begin{rem}
 A~Lie filtration is called regular if all the $\F^j$ correspond to the sections of subbundles $H^j \subset TM$. With a~regular Lie filtration, the osculating groups all have the same dimension as $M$ and form a~Lie groupoid over $M$. With a~singular Lie filtration, we can consider the open dense subset $\bigcap_{j\geq 1} M_{\F^j}$ formed by the intersection of the respective regular sets. On this subset, the singular Lie filtration is a~regular Lie filtration.
\end{rem}

Now to define the Helffer--Nourrigat cone, consider the locally compact space $\adiab\F^*$. We define~\smash{$\T^*_x\F := \overline{T^*M\times \R^*_+} \cap \adiab\F^*_{(x,0)}$}
and identify it to a~subset of $\gr_x(\Fb)^*$. By~\cite[Proposition~1.12]{AMYMaximalHypoellipticity}, this set is stable under the coadjoint action so we can consider the corresponding subset of representations through Kirillov theory.

\begin{mydef}\label{Def: ConeHF Intrinsic}
 The Helffer--Nourrigat cone is the set
 \[\HN(\Fb) := {\T^*\F}/{\Ad^*(\gr(\Fb))} \subset \widehat{\gr(\Fb)}.\]
 We also define its subset
 \[\HN_0(\Fb) := {(\T^*\F\setminus (M\times \{0\}))}/{\Ad^*(\gr(\Fb))},\]
 which corresponds to removing the trivial representation of each osculating group from the Helffer--Nourrigat cone.
\end{mydef}

\begin{rem}\label{Remark : One cannot use all the representations}
 Consider the simple case where $M = \R$ and define a~filtration of depth 3 by defining $\F^1$, $\F^2$, $\F^3$ as generated by $x^2\partial_x$, $x\partial_x$, $\partial_x$, respectively. The singular part is $\{0\}$. The osculating group at that point is the abelian group $\R^3$ with each component having grading~1,~2 and~3 (generated by the respective equivalence classes of $x^2\partial_x$, $x\partial_x$, $\partial_x$).

 We get a~differential operator of order 4 by setting $D= \bigl(x^2\partial_x\bigr)\partial_x - (x\partial_x)^2$. One can then ``freeze the coefficients'' at $x = 0$. The osculating Lie algebra is generated by the equivalence classes $\big[x^2\partial_x\big]$, $[x\partial_x]$, $[\partial_x]$. The operator $D$ yields the left-invariant operator on the osculating group $\big[x^2\partial_x\big][\partial_x] - [x\partial_x]^2$. Since the osculating group is abelian, we can take the Fourier transform of this operator. We obtain the function $\xi \mapsto \xi_1\xi_3 - \xi_2^2$ in the dual coordinates. It turns out however, after a~simple computation, that $D = -x\partial_x$ and has thus order~2. Its principal symbol, as an operator of order~4, would then have to vanish. Therefore, if we want a~reasonable notion of principal symbol, we can only take representations (which here corresponds to evaluations of the Fourier transform above) corresponding to $\xi \in \R^3$ with $\xi_1\xi_3 - \xi_2^2 = 0$. This turns out to be the Helffer--Nourrigat cone (at $x=0$) in that particular case.
\end{rem}

Since the definition of $\T^*\F$ is invariant under the zoom action, we get the following.

\begin{prop}
 The set $\T^*\F$ is invariant under $\delta^*_{\lambda}$, $\lambda>0$. Consequently, so are the sets $\HN(\Fb)$, $\HN_0(\Fb)$. The action on the quotient is still denoted by $\delta$.
\end{prop}
\begin{proof}
 The adiabatic foliation is preserved by the zoom action, and so is its dual with the dual action. Therefore, at $t=0$, we get the $\R^*_+$-action $\delta_{\lambda}, \lambda > 0$ on $\gr(\Fb)$, and both $\T^*\F$ and $\HN(\Fb)$ are stable under the dual action.
\end{proof}

\subsection{The Rothschild--Stein method}\label{Subsection:Rothschild--Stein}

We now give an alternate construction of the Helffer--Nourrigat cone based on the work of Rothschild and Stein~\cite{RothschildStein} that corresponds to the initial definition of Helffer and Nourrigat~\cite{HelfferNourrigatPaper}. The idea is to choose generators of each $\F^j$ to construct a~bigger graded Lie group that surjects onto every osculating group. Dualizing this projection, we obtain an embedding of the unitary representations of each osculating group into the representations of this bigger group. We can then give an alternate description of the Helffer--Nourrigat cone, where the group does not depend on the point in $M$. We now assume the modules of vector fields are (globally) finitely generated and use the following notations: let $N$ denote the depth of the filtration, i.e., $\F^j = \mathfrak{X}(M)$, $\forall j \geq N$. Now consider as in Example~\ref{Example : Finitely generated} vector fields $X_1,\dots,X_m \in \mathfrak{X}(M)$ and weights $\nu_1,\dots,\nu_m \in \N$ such that the submodule $\F^k$ is generated by the $X_I$ of weight at most~$k$, $I \subset \{1,\dots,m\}^n$, $n\in \N$. Now let $\mathfrak{g}$ be the universal graded Lie algebra of depth $N$ generated by elements $Y_k$, $1\leq k \leq m $ with $Y_k$ having the same degree as~$X_k$, i.e.,~$\nu_k$.

\begin{rem}
 In the original approach of~\cite{RothschildStein}, they only choose a~generating family for $\F^1$. This is because they consider the sub-Riemannian situation where $\F^{j+1} = \F^j + [\F^j,\F^1]$, $\forall j \geq 1$, i.e., all the groups appearing are stratified Lie groups.

 It is also worth mentioning that we only need a~generating family for the vector fields. In particular, given two different generating families, one can consider the union of both of these families. This gives rise to 3 universal nilpotent Lie algebras, $\g_1$, $\g_2$, and $\g$, respectively, and $\g$ projects on both $\g_1$ and $\g_2$. It is not the algebraic sum though, as it would (in general) not be a~universal Lie algebra.
\end{rem}

We define a~linear map $\beta \colon \g \to \mathfrak{X}(M)$ by sending $Y_k$ to $X_k$ and every commutator of weight at most $N$ onto the respective commutator of $X_j$'s. This map is not a~Lie algebra morphism because the Lie bracket of vector fields on $M$ with length bigger than $N$ will not necessarily vanish, whereas it automatically does in $\g$.
For $x \in M$, we get a~map $\beta_x := \ev_x \circ \beta \colon \g \to T_xM$. This map is surjective, so by duality we get a~linear injection $\beta_x^* \colon T^*_xM \hookrightarrow \g^*$.

\begin{mydef}\label{Def: ConeHN Extrinsic}
 The (extrinsic) Helffer--Nourrigat cone, at a~point $x\in M$ is the set of elements~${\ell \in \g^*}$ such that there exists a~sequence \smash{$((t_n,x_n,\xi_n))_{n\in \N} \in \bigl(\R^*_+ \times T^*M\bigr)^{\N}$} with $t_n\rightarrow 0$, $x_n\rightarrow x$ and $\delta_{t_n}^*\beta_{x_n}^*(\xi_n) \rightarrow \ell$.
 This set is stable under coadjoint action on $\g^*$ so we see it as a~subset $\HN_{\g}(\Fb)_x \subset \widehat{\g}$. Similarly as in Definition~\ref{Def: ConeHF Intrinsic}, we also define its subset $\HN_{0,\g}(\Fb)_x$, and denote by~$\HN_{\g}(\F)$ and $\HN_{0,\g}(\F)$ the disjoint union of the corresponding sets for $x\in M$.
\end{mydef}

In this definition, $\delta$ denotes the natural $\R^*_+$-action on the graded group $\g$, similarly as for the osculating groups.

\begin{rem}
 In general the only requirements on the group $\g$ should be that it is a~(finite-dimensional) graded Lie group of depth $N$, equipped with a~linear map $\g \to \mathfrak{X}(M)$ sending $\g_j$ to $\F^j$. This map should preserve commutators of length lower or equal to $N$. Finally, every localized map $\g \to T_xM$, $x\in M$ should be surjective.
\end{rem}

To define this version of the Helffer--Nourrigat cone, we had to make choices in order to construct the group $\g$. We now show that Definition~\ref{Def: ConeHN Extrinsic} coincides with Definition~\ref{Def: ConeHF Intrinsic} and thus does not depend on any choice. To do this, first notice that since $\beta$ sends $\g^j$ into $\F^j$ for each~${1\leq j \leq N}$ then it induces a~map $\gr(\beta) \colon \g \to \gr(\Fb)$.
This map sends each $Y_j$ to $X_j \mod \F^{k_j-1}$ where $1\leq k_j\leq N$ is the smallest integer for which $X_j \in \F^{k_j}$. We can also localize this map to get a~collection of maps
\[\gr(\beta)_x \colon\ \g \to \gr(\Fb)_x,\qquad x \in M.\]
All these maps are surjective by construction. Moreover, since $\beta$ preserves Lie brackets of length lower or equal to $N$ then each $\gr(\beta)_x$ is a~Lie algebra homomorphism. Indeed the Lie brackets of length lower or equal to $N$ are preserved, and the ones of higher length are equal to $0$ on both sides. As such, $\gr(\beta)_x$ also induces a~group homomorphism between the respective (connected, simply connected) Lie groups, denoted the same way.

Either via precomposition, or via the transpose map $\tsp{}\gr(\beta)_x $ and the use of the orbit method, the map $\gr(\beta)_x$ induces a~continuous map
\[\widehat{\gr(\beta)_x} \colon\ \widehat{\gr(\Fb)_x} \to \widehat{\g},\qquad x \in M.\]

\begin{prop}\label{Proposition : Independence HN}
 Let $x\in M$, the map \smash{$\widehat{\gr(\beta)_x}$} is a~topological embedding. Moreover, it induces a~homeomorphism between $\HN(\Fb)_x$ and $\HN_{\g}(\Fb)_x$. In particular, the latter does not depend on the choice of generators of $\Fb$ {\rm(}up to homeomorphism$)$. This homeomorphisms also identifies the subsets $\HN_0(\F)$ and $\HN_{0,\g}(\F)$.
\end{prop}

\begin{proof}
 We use the orbit method. Since $\gr(\beta)_x$ is surjective, then $\tsp\gr(\beta)_x\colon \gr_x(\Fb)^* \to \g^*$ is injective. Then if $\xi \in \gr_x(\Fb)^*$, $g \in \g$,
 \[\tsp\gr(\beta)_x\bigl(\Ad^*_{\gr(\beta)_x(g)}\xi\bigr) = \Ad^*_g(\tsp\gr(\beta)_x(\xi)).
 \]
 Therefore, two points $\xi,\xi' \in \gr_x(\Fb)^*$ lie in the same coadjoint orbit if and only if the points $\tsp\gr(\beta)_x(\xi),\tsp\gr(\beta)_x(\xi') \in \g^*$ lie in the same coadjoint orbit. Consequently, the corresponding quotient map \smash{$\widehat{\gr(\beta)} \colon \widehat{\gr_x(\Fb)} \to \widehat{\g}$} is injective and is also a~topological embedding.

 To show that the Helffer--Nourrigat cones are the same, we extend $\beta$ to $\widetilde{\beta} \colon M \times \g \times \R \to TM \times \R$ by $\widetilde{\beta}(x,X,t) = (x,\beta_x(\delta_t(X)),t)$.
 The map $\widetilde{\beta}$ takes values (when applied to sections) in the adiabatic foliation $\mathfrak{a}\F$. The map $\widetilde{\beta}^* \colon \mathfrak{a}\F^*\to M \times \g^* \times \R$ is also an embedding. At~${t=0}$, we obtain the previous map $\tsp\gr(\beta)$. At~${t\neq 0}$, we obtain the map $\delta^*_t\beta^*$. Therefore, since the limit conditions defining $\HN_{\g}(\F)$ and~$\HN(\F)$ are the same, we get
 \[\widehat{\gr(\beta)}(\HN(\Fb)) = \HN_{\g}(\Fb).\tag*{\qed}
 \]\renewcommand{\qed}{}
\end{proof}

\begin{rem}
 The independence of the choice of generators can also be proved without the intrinsic version of the Helffer--Nourrigat cones. If we get two sets of generators, we can construct the respective groups $\g_1$, $\g_2$. We can also create another nilpotent group $\g$ by considering the universal Lie algebra generated by both the generators used for $\g_1$ and $\g_2$. Given some $x \in M$, the map $\g \to \gr(\Fb)_x$ factors through both the maps $\g_i \to \gr(\Fb)_x$. Dually, we get a~commutative diagram
 \[\xymatrix{\widehat{\g} & \ar[l]_{\iota^1_x} \widehat{\g_1} \\
 \widehat{\g_2}\ar[u]^{\iota^2_x} & \ar[u] \ar[l]\widehat{\gr(\Fb)_x}.}\]
 The maps $\iota^i_x$ are both topological embeddings and, even though the images of $\iota_x^1$ and $\iota_x^2$ might differ, we have
 \[\iota_x^1(\HN_{(0),\g_1}(\Fb)) = \iota_x^2(\HN_{(0),\g_2}(\Fb)).\]
 Indeed, both are equal to the image of \smash{$\HN_{(0)}(\Fb)_x$} under the map \smash{$\widehat{\gr(\Fb)_x} \to \widehat{\g}$}.
\end{rem}

From now on, we will omit the subscript $\g$ in the Helffer--Nourrigat cone and will not make a~distinction between its intrinsic and extrinsic definition.

\section{Pedersen stratification of a~group}\label{Section:PedersenStratification}

Let $\g$ be a~graded Lie algebra. Recall that a~Jordan--H\"older basis is a~basis $Y_1, \dots, Y_n$ of $\g$ seen as a~vector space, such that every $\g_j = \lag Y_1,\dots,Y_j\rag$ is an ideal of $\g$. The idea of the Pedersen stratification is to regroup elements $\xi \in \g^*$ based on the numbers
\[d_{j,i}(\xi) := \dim\bigl(\Ad^*(\g_i) \bigl(\xi \mod \g_j^{\perp} \bigr)\bigr).
\]
The subspaces $\g_j$ being ideals, they are stable under the adjoint action. Therefore, $\g_j^{\perp}\subset\g$ is stable under the coadjoint action. Therefore, we get an action of $\g$ (hence of each $\g_i$) on the quotient \smash{${\g^*}/{\g^{\perp}_j}$}. We can see $\g_j^*\subset \g^*$ because we have chosen a~basis of $\g$, we then have~\smash{$\g_j^* \cong {\g^*}/{\g^{\perp}_j}$} and $d_{j,i}$ is the dimension of the orbit under $\g_i$ of the projection of $\xi$ onto $\g_j^*$.

The following properties of these numbers can be found in~\cite{Bonnet,PedersenStratification}.

\begin{prop}
 Let $\xi \in \g^*$, the numbers $d_{j,i} = d_{j,i}(\xi)$, $1\leq i,j\leq n$, satisfy
 \begin{enumerate}\itemsep=0pt
 \item[$(1)$] $d_{i,j} = d_{j,i}$.
 \item[$(2)$] $d_{j,i} - d_{j-1,i} \in \{0;1\}$.
 \item[$(3)$] $d_{j,j} - d_{j-1,j-1} \in \{0;2\}$.
 \item[$(4)$] If $d_{j,i} = d_{j-1,i}$, then $d_{j,k} = d_{j-1,k}$ for all $1\leq k \leq i$.
 \item[$(5)$] For $1\leq i \leq n$, $d_{i,i} = d_{i-1,i}+1$ if and only if $d_{i,i} = d_{i-1,i-1}+2$.
 \item[$(6)$] For all $g \in \g$, $d_{j,i}(\Ad^*_g(\xi)) = d_{j,i}(\xi)$.
 \end{enumerate}
\end{prop}

Before proving the result, notice that the last property implies that the numbers $d_{j,i}(\xi)$ are invariants of the coadjoint orbit corresponding to $\xi$. The idea of Pedersen's stratification is to regroup all the points (hence all the orbits) in $\g^*$ that have the same $d_{j,i}$. However, we also want to order them to get not just a~union of orbits but a~stratification.

\begin{proof}
 We can compute the dimension of the orbit by computing the dimension of its tangent space. Let $M(\xi)$ be the matrix in the Jordan--H\"older basis of the bilinear map
 $(X,Y) \mapsto \lag\xi,[X,Y]\rag$.
 Let $M_{(j,i)}$ be the sub-matrix of $M(\xi)$ formed by the $j$-th first rows and $i$-th first columns. Then $d_{j,i} = \rk(M_{(j,i)}(\xi))$. Since $M(\xi)$ is antisymmetric, this gives (1).

 The number $d_{j,i}-d_{j-1,i}$ is the dimension of the orbit under $\g_i$ of the image of $\xi$ in \smash{${\g_j^*}/{\g_{j-1}^*}$}.~This space being of dimension 1, the dimension of the orbit is either $0$ or $1$. This proves item~(2).

 The matrix $M(\xi)$ being antisymmetric, the rank $d_{j,j}(\xi)$ of the sub-matrices $M_{(j,j)}(\xi)$ (still antisymmetric) can only increase by $2$ (or not change at all) when we go from $j-1$ to $j$. This proves (3).

 If $d_{j,i} = d_{j-1,i}$, this means that the projection of $\xi$ in ${\g_j^*}/{\g_{j-1}^*}$ is a~fixed point for the action of~$\g_i$ (the orbit has dimension $0$). Then it is also the same for the subgroups $\g_k$, $1\leq k \leq i$. This~proves~(4).

 The matrix $M_{(i,i-1)}$ is an intermediate sub-matrix between $M_{(i-1,i-1)}$ and $M_{(i,i)}$ so (5) follows directly from the previous points by considering the contrapositions.

 Finally, let $g\in \g$, we have
 \begin{align*}
 \Ad^*(\g_i)\Ad^*_g(\xi) &= \Ad^*_g\Ad^*_{g^{-1}}\Ad^*(\g_i)\Ad^*_g(\xi) = \Ad^*(g)\Ad^*\bigl(g\g_ig^{-1}\bigr)\xi= \Ad^*(g)\Ad^*(\g_i)\xi.
 \end{align*}
 Since $\Ad^*_g$ is a~diffeomorphism of $\g^*$, it preserves the dimension and we get (6).
\end{proof}

\begin{ex}
 Let $\g = \g_1\oplus\cdots\oplus \g_N$ be a~graded Lie algebra. Let $Y_1, \dots, Y_n$ be a~basis of~$\g$ of eigenvectors for the inhomogeneous dilations. We have $\delta_{\lambda}(Y_i) = \lambda^{\nu_i}Y_i$, $\lambda>0$, $1\leq 1 \leq n$. If the $(\nu_i)_{1\leq i \leq n}$ form a~non-increasing sequence then $(Y_i)_{1\leq i \leq n}$ is a~Jordan--H\"older basis. With a~graded Lie algebra, the weights $\nu_i$ are integers and correspond to the grading of the elements~$Y_i$, i.e., $Y_i \in \g_{\nu_i}$. The way to obtain such a~basis is to first take a~basis of $\g_N$, then complete it with a~basis of $\g_{N-1}$, and so on until we have a~basis of $\g$.

 We call such a~Jordan--H\"older basis compatible with the grading.
\end{ex}

We now consider the indices for which a~jump in dimension happens. We define
\[J_i(\xi) = \{1\leq j \leq i,\, d_{j,i} = d_{j-1,i} +1 \},\qquad 1\leq i \leq n.\]
We can reconstruct the numbers $d_{j,i}(\xi)$ from these sets using properties (1)--(5). We put a~partial order on subsets $I, I' \subset \{1;\dots;n\}$ by $I\prec I' \Leftrightarrow \min(I\setminus I') \leq \min(I' \setminus I)$.
With the convention~${\min(\varnothing) = +\infty}$, this defines a~total ordering for which the maximal element is $\varnothing$.

Now the collection of all our orbits form a~set $\mathcal{J}(\xi)= (J_n(\xi),\dots,J_1(\xi)) \in \mathcal{P}(\{1,\dots,n\})^n$.

Notice that using (6), we have that for all $\xi \in \g^*$ and $g\in \g$, $\mathcal{J}(\Ad^*_g(\xi)) = \mathcal{J}(\xi)$.

We put the lexicographic order on $\mathcal{P}(\{1,\dots,n\})^n$ using the previous order $\prec$ on each component.

\begin{mydef}[Pedersen stratification]
 Let $\g$ be a~nilpotent group with a~choice of Jordan--H\"older basis. For $\mathcal{J}_0 \in \mathcal{P}(\{1,\dots,n\})^n$ we set $\Lambda_{\mathcal{J}_0} = \{\xi \in \g^* \mid \mathcal{J}(\xi) = \mathcal{J}_0\}$. The non-empty sets obtained that way can be ordered using the previous ordering for the indexing $\mathcal{J}_0$. We denote by $\Lambda_1, \dots, \Lambda_d$,
 the subsets obtained and ordered that way. We have $\g^* = \bigsqcup_{k = 1}^d \Lambda_k$. Since these sets are stable under the coadjoint action, we also define
 \[\Omega_k =  {\Lambda_k}/{\Ad^*(\g)},\qquad 1\leq k \leq d.\]
 We call the decomposition $\widehat{\g} = \bigsqcup_{k = 1}^d \Omega_k$ the Pedersen stratification of $\widehat{\g}$.
\end{mydef}

\begin{ex}
 Regardless of the Jordan--H\"older basis, we have $\Lambda_{(\varnothing,\dots,\varnothing)} = \Omega_{(\varnothing,\dots,\varnothing)} = [\g,\g]^{\perp}$.
 Indeed, the points on both sides correspond to the fixed point for the coadjoint action. The set~${(\varnothing,\dots,\varnothing)}$ is necessarily the biggest for the previous order.
\end{ex}

\begin{rem}
 Although the final stratum in the stratification does not depend on the choice of the Jordan--H\"older basis, the others do. For instance, the first stratum is always a~subset of the flat orbits (see~\cite{GoffengKuzmin}) but might depend on the Jordan--H\"older basis, see Example~\ref{Ex: Complex Heisenberg group} below.
\end{rem}

\begin{thm}[Pedersen~\cite{PedersenStratification}]\label{Theorem:ClassicPedersecStrat}
 Let $\g$ be a~nilpotent group with a~choice of Jordan--H\"older basis. Let $\Lambda_k$, $\Omega_k$, $1\leq k \leq d$ denote the Pedersen strata, in $\g^*$ and $\widehat{\g}$ respectively. We have:
 \begin{itemize}
\itemsep=0pt
 \item The sets \smash{$W_k = \bigcup_{\ell\leq k} \Lambda_{\ell}$} are Zariski open in $\g^*$. Consequently, the sets $U_k = \smash{{W_k}/{\Ad^*(\g)}} = \smash{ \bigsqcup_{\ell\leq k}\Omega_{\ell}}$ are open in $\widehat{\g}$.
 \item For every $1\leq k \leq d$, there is a~decomposition $\g^* = V_k^S\oplus V_k^T$ such that each orbit in $\Lambda_k$ meets $V_k^T$ exactly once. The corresponding set $C_k = \Lambda_k \cap V_k^T$ is then a~cross-section for the quotient map $\Lambda_k \to \Omega_k$.
 \item There is a~non-singular, birational bijection $\Psi_k \colon C_k \times V_k^S \to \Lambda_k$ such that for each $v\in C_k$, the map $p_k^T \circ \Psi_k(v,\cdot) \colon V_k^S \to V_k^T$,
 is a~polynomial, whose graph is the orbit $\g \cdot v$. Here $p_k^T$ denotes the projection onto $V_k^T$ parallel to $V_k^S$. The map $\pr_1\circ \Psi_k^{-1} \colon \Lambda_k \to C_k$
 is $\Ad^*(\g)$-invariant and the quotient is the inverse bijection of $C_k \isomto \Omega_k$.
 \end{itemize}
\end{thm}

This theorem shows that we can identify the sets $\Omega_k\subset \widehat{\g}$ to a~semi-algebraic set ${C_k \subset \g^*}$ and that, after some birational reparameterization, the set of corresponding orbits $\Lambda_k$ is homeo\-morphic to $C_k \times V_k^S$. This reparameterization is important as the orbits don't have to be flat. The first property is also crucial in what will follow as it will enable us to define ideals in the group $C^*$-algebra. It also gives sense to the term \emph{stratification} since we are adding ``layers'' to an increasing union of open sets in $\widehat{\g}$.

\begin{rem}
 There is another stratification of the orbits due to Puk\'anszky~\cite{Pukanszky}. It is obtained by regrouping points (hence orbits) $\xi \in\g^*$ that share the same set $J_n(\xi)$ and the order is the same as before. Therefore, Pedersen's stratification is finer than the one of Puk\'anszky in the sense that his strata are subsets of strata in the Puk\'anszky's stratification. The former is often called the \emph{fine} stratification while the latter is the \emph{coarse} stratification. The coarse stratification has similar properties as the fine one, the previous theorem is true for the coarse stratification as well. In most simple examples, the two coincide. See Example~\ref{Ex: Complex Heisenberg group} below to witness a~difference.
\end{rem}

The advantage of using the fine stratification instead of the coarse one can be seen through the next result.

\begin{lem}
 Let $\g$ be a~graded nilpotent group and fix a~Jordan--H\"older basis compatible with the grading. The inhomogeneous dilations $\delta_{\lambda} \in \Aut(\g)$, $\lambda > 0$, define an action of the positive real numbers on $\widehat{\g}$ by precomposition. The Kirillov homeomorphism \smash{$\widehat{\g} \cong {\g^*}/{\Ad^*(G)}$} is equivariant with the action of $\bigl(\tsp\delta_{\lambda}\bigr)_{\lambda > 0}$ on $\g^*$.
 Moreover, the sets $\Lambda_k$, $\Omega_k$, $W_k$, $U_k$ from Theorem~{\rm\ref{Theorem:ClassicPedersecStrat}} are all invariant under the respective $\R^*_+$-action.
 Finally, the action $\R^*_+ \act \Omega_k$ is free and proper for~${k < d}$ and so is the action $\R^*_+\act \Omega_d\setminus\{0\}$.
\end{lem}

\begin{proof}
 The equivariance of Kirillov's map is clear from its construction. Let $\xi \in \g^*$, let $\n \subset \g$ be a~subalgebra for which $\lag\xi,[\n,\n]\rag = 0$ and maximal for that property (we call $\n$ a~polarization for the orbit). Then $\xi$ defines a~character $\chi_{\xi,\n}$ of $\n$ by $X \mapsto \exp(2\ii \pi \xi(X))$. The representation corresponding to the orbit $\Ad^*(\g)\cdot \xi$ is then the induced representation from $\n$ to $\g$: $\mathrm{Ind}_{\n}^{\g}(\chi_{\xi,\n})$. This representation can be shown (cf.\ \cite{Kirillov}) to be irreducible, and not depend on the choice of $\xi$ or $\n$, up to unitary equivalence. Here if we fix $\lambda > 0$ and take $\xi$ and $\n$, then the algebra $\delta_{\lambda^{-1}}(\n)$ is a~polarization for $\tsp\delta_{\lambda}(\xi)$. Moreover, $\chi_{\tsp\delta_{\lambda}(\xi),\delta_{\lambda^{-1}}(\n)} = \chi_{\xi,\n}\circ \delta_{\lambda}$.
 We conclude by functoriality of the induction.

 Since the Jordan--H\"older basis is compatible with the grading, we have $\forall 1\leq i, j\leq n$, $ d_{j,i}\bigl(\tsp\delta_{\lambda}(\xi)\bigr) = d_{j,i}(\xi)$,
 and the rest follows.
\end{proof}

























































































\section{Pedersen stratification of the Helffer--Nourrigat cone}\label{Section : Stratification HN Cone}

Let $(M,\Fb)$ be a~singular Lie filtration. According to Section~\ref{Section: Helffer--Nourrigat cone}, the Helffer--Nourrigat cone is defined as a~closed subset $\HN(\Fb)$ of the coadjoint orbits of $\gr(\Fb)$. As before, we make the assumption that $\Fb$ is globally finitely generated so we may use the Rothschild--Stein method to find a~global graded Lie basis with model group $\g$. We may then take a~Pedersen stratification of $\widehat{\g} = \bigsqcup_{j=1}^d\Omega_j^{\g}$. Since $\HN(\Fb) \subset M\times \widehat{\g}$ is closed, then all the sets
\[
\Omega_j := \bigl(M\times\Omega_j^{\g}\bigr)\bigcap \HN(\Fb),
\]
are closed subsets of $M \times \Omega_j^{\g}$ hence locally compact and Hausdorff. Moreover, they form a~partition of $\HN(\Fb)$.

Notice however that since we use a~universal Lie algebra construction, it can be fairly big, so some of the Pedersen strata may not intersect $\HN(\Fb)$. We can solve this problem by getting rid of the empty strata and re-indexing the remaining ones. Another thing that may occur is that a~strata $\Omega_j$ could have an empty fiber over a~point $x \in M$ but be non-empty (so have non-empty fiber over other points). See Section~\ref{Section: ExamplesCones} below, for examples, where this phenomenon happens. Typically in a~lot of sub-Riemannian settings, $\F^1$ generates the whole tangent space by itself on an open dense subset $M_{\rm reg}\subset M$ and we only need to consider Lie brackets over the singular set. In that case, the only non-empty stratum over $M_{\rm reg}$ is the last one since it only consists of abelian representations. However, we could have infinite-dimensional representations appearing in addition over the singular set.

The goal of this section is to prove Theorem~\ref{Theorem:Stratification}. The sets $\Omega_j$ defined above define a~stratification of the Helffer--Nourrigat cone. However, given a~point $x \in M$, we want to relate the fibers over $x$ to the Pedersen strata for $\gr(\Fb)_x$. This is a~consequence of the following result.

\begin{lem}
 Let $\g,\mathfrak{h}$ be graded Lie algebras, $\varphi\colon \g\twoheadrightarrow\mathfrak{h}$ a~surjective homomorphism that preserves the grading. Let $Y_1, \dots, Y_N$ be a~Jordan--H\"older basis of $\g$ compatible with the grading. Then there exists a~Jordan--H\"older basis of $\mathfrak{h}$ such that $\widehat{\varphi} \colon \widehat{\mathfrak{h}} \hookrightarrow \widehat{\g}$ maps elements belonging to a~same Pedersen stratum in $\widehat{\mathfrak{h}}$ to the same Pedersen stratum in $\widehat{\g}$. Moreover, the induced map between the Pedersen strata of $\widehat{\mathfrak{h}}$ and $\widehat{\g}$ is injective.
\end{lem}

\begin{proof}
 Define $Z_j = \varphi(Y_j)$, $1\leq j \leq N$. These vectors span the Lie algebra $\mathfrak{h}$. We extract a~basis in the following way. Define $i_1 = \min\{i\mid Z_i \neq 0\}$,
 and then inductively $i_{k+1} = \min\{i \mid Z_i \notin \Span(Z_{i_1},\dots, Z_{i_k})\}$,
 until we get a~basis for some $i_n$.
 The vectors $Z_{i_1}, \dots, Z_{i_n}$ form a~Jordan--H\"older basis of $\mathfrak{h}$, compatible with the grading.

 We now want to relate jump indices for the two bases when we apply $\tsp\varphi$. To do that, we define
$u \colon  \{1,\dots,N\}  \to \{1,\dots,n\}$, $ i \mapsto \max\{k\mid i_k \leq i\}$.
 The map $u$ is a~right inverse to $\sigma \colon k \mapsto i_k$. We have $\varphi(\g_i) = \mathfrak{h}_{u(i)}$ for all $1\leq i \leq N$. Recall also that if $\xi \in \mathfrak{h}^*$ and $g \in \g$ then \smash{$\tsp\varphi\bigl(\Ad^*_{\varphi(g)}(\xi)\bigr) = \Ad^*_g\bigl(\tsp\varphi(\xi)\bigr)$},
 i.e., $\tsp\varphi$ induces a~homeomorphism between the $\mathfrak{h}$-coadjoint orbit of $\xi$ and the $\g$-coadjoint orbit of~$\tsp\varphi(\xi)$.
 Consequently, for all $\xi \in \mathfrak{h}^*$ and $1\leq i \leq N$ we have $J^i\bigl(\tsp(\varphi(\xi)\bigr) = \sigma\bigl(J^{u(i)}(\xi)\bigr)$.
 We therefore map Pedersen strata to Pedersen strata, injectively between the strata.
\end{proof}

\begin{rem}
 The proof of the lemma is for the Pedersen stratification, i.e., the fine stratification. One can see that the same proof also works for the Puk\'anszky stratification, i.e., the coarse stratification.
\end{rem}

Using this result we have proved the following, which is a~reformulation of Theorem~\ref{Theorem:Stratification}.

\begin{thm}\label{Theorem:Stratification-InText}
 Let $(M,\Fb)$ be a~finitely generated singular filtration. We fix generators and let~$\g$ be the group obtained by the Rothschild--Stein method.
 Let $\Omega_j \subset \HN(\Fb)$, $1\leq j \leq d$, be the non-empty sets obtained by intersecting $\HN(\Fb)$ with the Pedersen strata. Then the sets $U_j=\bigcup_{k\leq j}\Omega_k$ are open and $U_d = \HN(\Fb)$, i.e., the sets $\Omega_j$ form a~partition of $\HN(\Fb)$.

 Moreover, if $x \in M_{\rm reg}$, the intersections $\Omega_j\cap \HN(\F)_x$ form a~Pedersen stratification of $\gr(\Fb)_x$. If $x\in M\setminus M_{\rm reg}$ then the intersections $\Omega_j\cap \HN(\F)_x$ are closed subsets of a~Pedersen stratification of $\gr(\Fb)_x$.
\end{thm}







































































































































































































































































































































































\begin{thebibliography}{99}
\footnotesize\itemsep=0pt

\bibitem{AMYMaximalHypoellipticity}
Androulidakis I., Mohsen O., Yuncken R., A~pseudodifferential calculus for
 maximally hypoelliptic operators and the {H}elffer--{N}ourrigat conjecture,
 \href{https://arxiv.org/abs/2201.12060}{arXiv:2201.12060}.

\bibitem{AndroulidakisSkandalisHolonomy}
Androulidakis I., Skandalis G., The holonomy groupoid of a~singular foliation,
 \href{https://doi.org/10.1515/CRELLE.2009.001}{\textit{J.~Reine Angew.
 Math.}} \textbf{626} (2009), 1--37,
 \href{https://arxiv.org/abs/math/0612370}{arXiv:math/0612370}.

\bibitem{AndroulidakisSkandalisPseudoDiff}
Androulidakis I., Skandalis G., Pseudodifferential calculus on a~singular
 foliation, \href{https://doi.org/10.4171/JNCG/72}{\textit{J.~Noncommut.
 Geom.}} \textbf{5} (2011), 125--152,
 \href{https://arxiv.org/abs/0909.1342}{arXiv:0909.1342}.

\bibitem{benedetto}
Benedetto L., Quantum ergodicity for contact metric structures,
 \href{https://arxiv.org/abs/2507.18216}{arXiv:2507.18216}.

\bibitem{Bonnet}
Bonnet P., Param\'etrisation du dual d'une alg\`ebre de {L}ie nilpotente,
 \href{https://doi.org/10.5802/aif.1144}{\textit{Ann. Inst. Fourier
 (Grenoble)}} \textbf{38} (1988), 169--197.

\bibitem{CGGP}
Christ M., Geller D., G{\l}owacki P., Polin L., Pseudodifferential operators on
 groups with dilations,
 \href{https://doi.org/10.1215/S0012-7094-92-06802-5}{\textit{Duke Math.~J.}}
 \textbf{68} (1992), 31--65.

\bibitem{CrenFilteredCrossedProduct}
Cren C., Filtered calculus and crossed products by {$\mathbb R$}-actions,
 \href{https://doi.org/10.7900/jot}{\textit{J.~Operator Theory}} \textbf{94}
 (2025), 473--511, \href{https://arxiv.org/abs/2308.13532}{arXiv:2308.13532}.

\bibitem{DaveHallerBGG}
Dave S., Haller S., Graded hypoellipticity of {BGG} sequences,
 \href{https://doi.org/10.1007/s10455-022-09870-0}{\textit{Ann. Global Anal.
 Geom.}} \textbf{62} (2022), 721--789,
 \href{https://arxiv.org/abs/1705.01659}{arXiv:1705.01659}.

\bibitem{deGraafClassification}
de~Graaf W.A., Classification of 6-dimensional nilpotent {L}ie algebras over
 fields of characteristic not~2,
 \href{https://doi.org/10.1016/j.jalgebra.2006.08.006}{\textit{J.~Algebra}}
 \textbf{309} (2007), 640--653,
 \href{https://arxiv.org/abs/math/0511668}{arXiv:math/0511668}.

\bibitem{DebordBourbaki}
Debord C., Hypoellipticit\'e de polyn\^omes de champs de vecteurs et
 conjectures de {H}elffer et {N}ourrigat (d'apr\`es {I}.~{A}ndroulidakis, {O}.~{M}ohsen, {R}.~{Y}uncken), \href{https://doi.org/10.24033/ast.1261}{\textit{Ast\'erisque}} \textbf{462} (2025), 289--318,
 \href{https://arxiv.org/abs/2504.13508}{arXiv:2504.13508}.

\bibitem{DixmierDouady}
Dixmier J., Douady A., Champs continus d'espaces hilbertiens et de~{$C^{\ast}
 $}-alg\`ebres, \href{https://doi.org/10.24033/bsmf.1596}{\textit{Bull. Soc.
 Math. France}} \textbf{91} (1963), 227--284.

\bibitem{Ewert}
Ewert E., Pseudo-differential extension for graded nilpotent {L}ie groups,
 \href{https://doi.org/10.4171/dm/940}{\textit{Doc. Math.}} \textbf{28}
 (2023), 1323--1379,
 \href{https://arxiv.org/abs/2002.01875}{arXiv:2002.01875}.

\bibitem{FermanianFischerGradedLie}
Fermanian-Kammerer C., Fischer V., Defect measures on graded {L}ie groups,
 \href{https://doi.org/10.2422/2036-2145.201711_011}{\textit{Ann. Sc. Norm.
 Super. Pisa Cl. Sci.}} \textbf{21} (2020), 207--291,
 \href{https://arxiv.org/abs/2145.20171}{arXiv:2145.20171}.

\bibitem{FermanianLetrouit}
Fermanian-Kammerer C., Letrouit C., Observability and controllability for the
 {S}chr\"odinger equation on quotients of groups of {H}eisenberg type,
 \href{https://doi.org/10.5802/jep.176}{\textit{J.~\'Ec. Polytech. Math.}}
 \textbf{8} (2021), 1459--1513,
 \href{https://arxiv.org/abs/2009.13877}{arXiv:2009.13877}.

\bibitem{GoffengKuzmin}
Goffeng M., Kuzmin A., Index theory of hypoelliptic operators on {C}arnot
 manifolds, \href{https://arxiv.org/abs/2203.04717}{arXiv:2203.04717}.

\bibitem{HelfferNourrigatPaper}
Helffer B., Nourrigat J., Hypoellipticit\'e maximale pour des op\'erateurs
 polyn\^omes de champs de vecteurs,
 \href{https://doi.org/10.5802/jedp.206}{\textit{C.~R.~Acad. Sci. Paris
 S\'er.~A-B}} \textbf{289} (1979), A775--A778.

\bibitem{HelfferNourrigatBook}
Helffer B., Nourrigat J., Hypoellipticit\'e maximale pour des op\'erateurs
 polyn\^omes de champs de vecteurs, \textit{Progr. Math.}, Vol.~58,
 Birkh\"auser, Boston, MA, 1985.

\bibitem{Hormander}
H\"ormander L., Hypoelliptic second order differential equations,
 \href{https://doi.org/10.1007/BF02392081}{\textit{Acta Math.}} \textbf{119}
 (1967), 147--171.

\bibitem{Kirillov}
Kirillov A.A., Unitary representations of nilpotent~{L}ie groups,
 \href{https://doi.org/10.1070/RM1962v017n04ABEH004118}{\textit{Russ. Math.
 Surv.}} \textbf{17} (1962), 53--104.

\bibitem{Letrouit}
Letrouit C., Subelliptic wave equations are never observable,
 \href{https://doi.org/10.2140/apde.2023.16.643}{\textit{Anal. PDE}}
 \textbf{16} (2023), 643--678,
 \href{https://arxiv.org/abs/2002.01259}{arXiv:2002.01259}.

\bibitem{LipsmanRosenberg}
Lipsman R.L., Rosenberg J., The behavior of {F}ourier transforms for nilpotent
 {L}ie groups,
 \href{https://doi.org/10.1090/S0002-9947-96-01583-8}{\textit{Trans. Amer.
 Math. Soc.}} \textbf{348} (1996), 1031--1050.

\bibitem{MohsenBlowupFoliation}
Mohsen O., Blow-up groupoid of singular foliations,
 \href{https://arxiv.org/abs/2105.05201}{arXiv:2105.05201}.

\bibitem{PedersenStratification}
Pedersen N.V., Geometric quantization and the universal enveloping algebra of a~nilpotent {L}ie group, \href{https://doi.org/10.2307/2001294}{\textit{Trans.
 Amer. Math. Soc.}} \textbf{315} (1989), 511--563.

\bibitem{Pukanszky}
Pukanszky L., Unitary representations of solvable~{L}ie groups,
 \href{https://doi.org/10.24033/asens.1218}{\textit{Ann. Sci. \'Ecole Norm.
 Sup.}} \textbf{4} (1971), 457--608.

\bibitem{RothschildStein}
Rothschild L.P., Stein E.M., Hypoelliptic differential operators and nilpotent
 groups, \href{https://doi.org/10.1007/BF02392419}{\textit{Acta Math.}}
 \textbf{137} (1976), 247--320.

\bibitem{Takai}
Takai H., On a~duality for crossed products of {$C^{\ast} $}-algebras,
 \href{https://doi.org/10.1016/0022-1236(75)90004-x}{\textit{J.~Functional
 Analysis}} \textbf{19} (1975), 25--39.

\bibitem{Taylor}
Taylor M.E., Noncommutative microlocal analysis.~{I},
 \href{https://doi.org/10.1090/memo/0313}{\textit{Mem. Amer. Math. Soc.}}
 \textbf{52} (1984), iv+182~pages.

\end{thebibliography}


